\section*{Chapitre \ref{ch:qm:probability-at-speed} --- Probabilité à grande vitesse}

\begin{solution}{exo:qm:probability-at-speed:1}
$M_1 = 13 + 4 \times (364 - 13)/51 = 13 + 27.53 = 40.53$, contre 35 avant l'opération.
\end{solution}

\begin{solution}{exo:qm:probability-at-speed:2}
(a) Oui~: on sait à chaque instant si cela s'est déjà produit. (b) Non~: le plus bas n'est connu qu'à
la clôture. (c) Oui~: un temps déterministe est un \omterm{def:qm:probability-at-speed:stopping}{temps d'arrêt}. (d) Oui, à condition que la prévision soit
elle-même adaptée à l'information observée.
\end{solution}

\begin{solution}{exo:qm:probability-at-speed:3}
$b/(a+b) = 3/4$. La marche arrêtée est bornée, donc le théorème d'arrêt optionnel donne $\E[S_\tau] = 0$~:
$\tfrac34 \times 1 - \tfrac14 \times 3 = 0$. Gagner souvent se paie par des pertes plus lourdes
quand on perd.
\end{solution}

\begin{solution}{exo:qm:probability-at-speed:4}
Les révisions sont orthogonales et de variances égales, si bien que les 30 dernières minutes résolvent $30/390 =
7.7\%$ de la variance.
\end{solution}

\begin{solution}{exo:qm:probability-at-speed:5}
$\E[Z] = e^{-c^2/2}\E[e^{cX}] = 1$ et $Z > 0$, donc $\mathbb Q$ est une mesure de
probabilité équivalente à $\P$. Sa \omterm{def:qm:probability-at-speed:charfn}{fonction caractéristique} est $e^{\iu uc - u^2/2}$
(l'\cref{ex:qm:probability-at-speed:shift})~: $X \sim \mathcal N(c,1)$ sous $\mathbb Q$, donc
$\E^{\mathbb Q}[X] = c$ et $\mathbb Q(X > c) = \tfrac12$.
\end{solution}

\begin{solution}{exo:qm:probability-at-speed:6}
$\E_4[F_5] = \E_4[M_5] = M_4$, donc $\E[F_5 - F_4 \mid \mathcal F_4] = (1 - a)(M_4 - 35) = \frac{1 -
a}{a}(F_4 - 35)$~: la révision est prévisible, et la pente vaut $(1 - a)/a = 0.667$ pour $a = 0.6$.
La même pente apparaît sur la révision précédente $F_4 - F_3 = a(M_4 - M_3)$, car $M_4 - 35$ est une
somme de révisions orthogonales dont $M_4 - M_3$ est l'une.
\end{solution}

\begin{solution}{exo:qm:probability-at-speed:7}
0.661 avec une statistique $t$ de 45.9, contre 0.667 en théorie~; la pente du prévisionniste honnête vaut
$-0.005$.
\end{solution}

\begin{solution}{exo:qm:probability-at-speed:8}
Une moyenne nulle est nécessaire, mais non suffisante~: le prévisionniste de l'exemple des cartes, qui sous-réagit, a aussi des
révisions de moyenne nulle. Les révisions d'une \omterm{def:qm:probability-at-speed:martingale}{martingale} doivent être imprévisibles à partir de tout ce qui est connu
avant elles, à commencer par ses propres révisions passées, et le test consiste en une régression des
révisions sur l'information antérieure, avec des erreurs types qui tiennent compte du groupement de
la volatilité (chapitre~11) et une correction pour le grand nombre de signaux testés (chapitre~12).
\end{solution}

\begin{solution}{pb:qm:probability-at-speed:1}
\textbf{1.} $\Var(C) = 1.80^2 + 0.30^2 = 3.33$~; écart-type de 1.82~USD.
\textbf{2.} Avant l'annonce, $M_t = S_t$, le cours observé en $t$ (les accroissements restants
et $J$ sont de moyenne nulle)~; après elle, $M_t = S_t + J$.
\textbf{3.} C'est la \omterm{def:qm:probability-at-speed:doob}{martingale de Doob} $\E_t[C]$ d'une variable intégrable.
\textbf{4.} $J$.
\textbf{5.} $\P(|J| > 0.14) = 2(1 - \Phi(0.14/0.30)) = 64.1\%$.
\textbf{6.} Les révisions sont orthogonales (la \cref{prop:qm:probability-at-speed:orthogonal}), donc
leurs variances s'ajoutent, et $M_0$ est une constante.
\textbf{7.} $150/390 \times 3.24/3.33 = 37.4\%$.
\textbf{8.} $0.09/3.33 = 2.7\%$.
\textbf{9.} $(0.09 + 10/390 \times 3.24)/3.33 = 5.2\%$~: 2.7~\% à l'annonce et 2.5~\% dans les
dix dernières minutes.
\textbf{10.} 5.23~\% dans les 20\,000 jours simulés.
\textbf{11.} Non~: sa révision sur les dix dernières minutes contient $J/2$, que sa
révision de 15~h~50 prédit $J/2$.
\textbf{12.} Les deux révisions valent $J/2$ et $J/2 + $ dix minutes d'accroissements~; la
corrélation vaut $0.15^2/(0.15\sqrt{0.15^2 + 0.0831}) = 0.46$.
\textbf{13.} Pente $\Cov/\Var = 0.15^2/0.15^2 = 1$~; la simulation donne 0.98 avec une
statistique $t$ de 72 sur 20\,000 jours.
\textbf{14.} Avec $t = r\sqrt{n - 1}/\sqrt{1 - r^2}$, $t = 2$ impose $n = 4(1 - r^2)/r^2 + 1 =
16$ jours.
\textbf{15.} La transaction rapporte $C - (S + J/2)$ multiplié par le signe de $J$, dont la moyenne vaut $\E|J|/2 =
0.5 \times 0.30\sqrt{2/\pi} = 0.120$~: 12~cents par action.
\textbf{16.} Les dix dernières minutes portent alors $3.24 \times 30/410 = 0.237$, et la part vaut
$(0.09 + 0.237)/3.33 = 9.8\%$.
\textbf{17.} Non, si les deux écarts types évoluent avec le cours~: les parts sont des rapports de
variances.
\textbf{18.} Ce qui se résout après le dernier moment où l'on peut agir est un risque que l'on porte, non
une information que l'on peut utiliser~; un trader qui décide à 15~h~50 fait face à 5.2~\% de la variance de la journée,
et non à 100~\%.
\textbf{19.} Résultat nommé~: \emph{fair value de la clôture}~: dans le modèle, 5.2~\% de la
variance du cours de clôture se résout entre l'annonce de 15~h~50 et la clôture (2.7~\% par
l'annonce elle-même), et 9.8~\% si les dix dernières minutes sont trois fois plus volatiles.
\textbf{20.} Elle promet que les révisions ne peuvent pas être prédites à partir de l'information propre au
prévisionniste~: personne disposant de cette information ne peut donc trader avec profit contre elle~; elle
ne dit rien de l'exactitude.
\end{solution}

\begin{solution}{iq:qm:probability-at-speed:1}
Avec un temps et un crédit illimités, on s'arrête en $+1$ presque sûrement, mais le \omterm{def:qm:probability-at-speed:stopping}{temps d'arrêt} est
non borné et les pertes antérieures ne sont pas uniformément intégrables~; avec un crédit fini $b$
on gagne $+1$ avec probabilité $b/(b+1)$ et on perd $b$ sinon~: espérance nulle.
\iqlookfor{arrêt optionnel, et pourquoi ses hypothèses comptent.}
\end{solution}

\begin{solution}{iq:qm:probability-at-speed:2}
Un \omterm{def:qm:probability-at-speed:filtration}{processus adapté} et intégrable dont l'\omterm{def:qm:probability-at-speed:condexp}{espérance conditionnelle} de toute valeur future est sa
valeur présente. Le cours d'une action a un rendement espéré positif sous la probabilité historique,
donc c'est une \omterm{def:qm:probability-at-speed:martingale}{sous-martingale} sous cette probabilité~; actualisé par le compte monétaire, c'est une \omterm{def:qm:probability-at-speed:martingale}{martingale}
sous la probabilité risque-neutre (chapitre~5).
\iqlookfor{la définition et le rôle de la mesure de probabilité.}
\end{solution}

\begin{solution}{iq:qm:probability-at-speed:3}
Pour $j < k$, $\E[D_jD_k] = \E[D_j\E_{k-1}[D_k]] = 0$. Une prévision qui est une \omterm{def:qm:probability-at-speed:condexp}{espérance conditionnelle} honnête
n'a que des révisions que personne ne peut prédire à partir de ses révisions passées~; si une régression
de la révision du jour sur celle de la veille a une pente, la prévision est inefficace.
\iqlookfor{retirer ce qui est connu de l'espérance, et la conséquence testable.}
\end{solution}

\begin{solution}{iq:qm:probability-at-speed:4}
La seconde révision du concurrent est prévisible à partir de la première, donc sa prévision n'est pas une
\omterm{def:qm:probability-at-speed:martingale}{martingale}~: elle sous-réagit. Trader contre son cours périmé à 15~h~50, dans le sens
de l'annonce~; dans le modèle du chapitre, cela rapporte la moitié de la valeur absolue moyenne de l'annonce, soit 12~cents par action.
\iqlookfor{des révisions prévisibles signalent une prévision exploitable.}
\end{solution}

\begin{solution}{iq:qm:probability-at-speed:5}
$\E[X \mid Y]$ est le meilleur prédicteur parmi toutes les fonctions de $Y$~; la régression linéaire est la
meilleure parmi les fonctions affines. Elles coïncident lorsque $\E[X \mid Y]$ est affine, comme c'est le cas lorsque $(X, Y)$ est
conjointement gaussien.
\iqlookfor{projections sur des sous-espaces emboîtés~; le cas gaussien.}
\end{solution}

\begin{solution}{iq:qm:probability-at-speed:6}
Conserver chaque valeur publiée avec son horodatage~; chaque jour, régresser les révisions sur des
horizons fixes sur la révision précédente et sur d'autres informations disponibles à la date antérieure
(déséquilibre du carnet d'ordres, transactions récentes), avec des erreurs types robustes, et présenter les parts de variance
selon l'heure de la journée. Le test est trompeur si les horizons se chevauchent (erreurs autocorrélées), si la
volatilité est hétéroscédastique et que les erreurs types l'ignorent, si les horodatages sont ceux
de la publication plutôt que de l'information (biais d'anticipation), si de nombreux signaux sont testés et que
le meilleur est rapporté, et si la valeur finale est elle-même révisée.
\iqlookfor{régression des révisions sur l'information antérieure, et façons dont le test échoue.}
\end{solution}
